\SourceLeaf{352}nahe hundert Jahren \textsc{Hiorter} und \textsc{Celsius} bemerkt, und nachher vielfache Beobachtungen bestätigt. Es liess sich hieraus schliessen, dass dieselben Kräfte, welche die Erscheinung eines Nordlichts hervorbringen, zugleich auch auf die Magnetnadel wirken, und dass diese Wirkungen sich auf sehr bedeutende Entfernungen erstrecken müssen, da die Nordlichter in einem weiten Umkreise sichtbar zu sein pflegen. Einen noch grössern Begriff von der weiten Ausdehnung der Wirksamkeit jener räthselhaften Kräfte erhalten wir durch die Bemerkung von Hrn. \textsc{Arago}, dass oft an denselben Tagen, wo er in Paris starke Störungen des regelmässigen Ganges der Magnetnadel beobachtet hatte, an entfernten Orten Nordlichter gesehen waren, deren Sichtbarkeit sich über den Horizont von Paris nicht erhoben hatte.

Die Unregelmässigkeiten in den Äusserungen des Erdmagnetismus, deren häufiges Vorkommen besonders auch Hr. \textsc{von Humboldt} bei seinen zahlreichen Beobachtungen der täglichen und stündlichen Bewegungen der Magnetnadel wahrgenommen hatte, erhielten hiedurch ein eigenthümliches Interesse. Wenn gleich jene Bemerkungen durchaus nicht dazu berechtigten, alle unregelmässigen Bewegungen, als gleichzeitig mit Nordlichtern zu betrachten, und die Möglichkeit noch nicht ausschlossen, dass viele, vielleicht die meisten, nur von localen Ursachen herrührten, so liess sich doch kaum verkennen, dass nicht selten grosse und fernhin wirkende Naturkräfte dabei im Spiel sind, deren Kenntniss, wenn auch noch nicht in Beziehung auf ihre Quelle, sondern zunächst nur in Beziehung auf die Verhältnisse ihrer Wirksamkeit und Verbreitung, einen würdigen Gegenstand der Naturforschung darbietet.

Obenhin und auf gut Glück gemachte Wahrnehmungen können uns diesem Ziele nicht näher bringen: um es zu erreichen, müssen viele solche Erscheinungen im genauen Detail an vielen Orten gleichzeitig verfolgt und nach Zeit und Grösse scharf gemessen werden. Dazu sind aber vorgängige bestimmte Verabredungen zwischen solchen Beobachtern, denen angemessene Hülfsmittel zu Gebote stehen, wesentlich nothwendig.

Der berühmte Naturforscher, dem unsere Kenntniss des Erdmagnetismus so viele Bereicherung verdankt, hat auch hier zuerst die Bahn gebrochen. Hr. \textsc{von Humboldt} errichtete in Berlin gegen Ende des Jahrs 1828 für die magnetischen Beobachtungen ein eignes eisenfreies Häuschen, stellte darin einen von \textsc{Gambey} verfertigten Variationscompass auf, und verband sich mit andern Be\SourceLeaf{353}sitzern ähnlicher Apparate an mehrern zum Theil sehr entlegenen Orten zu regelmässigen an verabredeten Tagen auszuführenden Beobachtungen der magnetischen Variation. Es wurden acht Termine im Jahre, jeder zu 44 Stunden. festgesetzt, an denen die magnetische Abweichung von Stunde zu Stunde aufgezeichnet werden sollte: an einigen Orten beobachtete man in noch engern Zeitgrenzen, von halber zu halber Stunde, oder von zwanzig zu zwanzig Minuten. Die nähern Bestimmungen findet man im 19. Bande von \textsc{Poggendorff}'s Annalen der Physik S. 361, und ebendaselbst auch die Beobachtungen, welche dieser Verabredung gemäss an den Terminen im Jahr 1829 und 1830 in Berlin, Freiberg, Petersburg, Kasan, und Nicolajef angestellt sind, so wie graphische Darstellungen von dreien derselben.

In dem hiesigen magnetischen Observatorium, welches im Jahr 1833 erbaut wurde und dessen magnetischer Apparat eine von den früher angewandten gänzlich verschiedene Einrichtung hat, wurden diese Terminsbeobachtungen zum erstenmal am 20. und 21. März 1834 vollständig angestellt, wozu correspondirende bloss aus Berlin bekannt geworden sind: aber in Göttingen war von zehn zu zehn Minuten, in Berlin nur von Stunde zu Stunde beobachtet. Gleichwohl zeigten diese Berliner Aufzeichnungen mehrere ziemlich beträchtliche Bewegungen, die man in den Göttinger Beobachtungen wiederfand, während diese letztern in den Zwischenzeiten eine grosse Menge anderer Bewegungen zu erkennen gaben, welche natürlich in Berlin ganz ausfallen mussten. Die Frage, ob ein kleinerer oder grösserer Theil der in Göttingen wahrgenommenen Schwankungen bloss local gewesen sei, blieb daher noch ohne Entscheidung.

Allein schon der nächste Termin, am 4. und 5. Mai, führte eine solche Entscheidung herbei. Die Zwischenzeiten wurden noch enger genommen, nemlich von fünf zu fünf Minuten, wodurch die Resultate noch bedeutend schärfer ausgeprägt erschienen. Correspondirende Beobachtungen mit \textsc{Gambey}'schen Apparaten sind von diesem Termine, eben so wie von allen spätern, überall keine mehr bekannt geworden. Dagegen hatte Hr. \textsc{Sartorius}, der an den Beobachtungen vom Märztermine in Göttingen thätigen Antheil genommen, und sich für eine mehrjährige nach Italien zu unternehmende Reise mit einem dem Göttingischen ganz ähnlichen, nur in kleinern Dimensionen gearbeiteten Apparate versehen hatte, mit diesem den Maitermin in Waltershausen (einem Gute in Baiern, etwa 20 Meilen von Göttingen entfernt) sorgfältig und vollständig in engen Zeitintervallen beob\SourceLeaf{354}achtet. Hier zeigte sich nun eine wirklich überraschend grosse Uebereinstimmung nicht nur in der grössern, sondern auch fast in sämmtlichen kleinern in kurzen Zeitfristen wechselnden Schwankungen, so dass in der That gar nichts übrig blieb, was man localen Ursachen beizumessen befugt gewesen wäre.

Die drei folgenden Termine im Junius, August und September wurden in Göttingen ganz auf dieselbe Weise abgehalten, während die Anzahl der auswärtigen Theilnehmer mit ähnlichen oder gleichen Apparaten sich fortwährend vergrösserte. Hr. Prof. \textsc{Encke} hatte sich, nachdem er die hiesigen Einrichtungen durch eigne Ansicht kennen gelernt hatte, für Berlin provisorisch einen ähnlichen Apparat nach kleinern Dimensionen anfertigen lassen; Hr. \textsc{Sartorius} beobachtete mit dem seinigen in allen Terminen, wo die Umstände es verstatteten (im Junius in Frankfurt, im September in Bramberg im Salzburgschen); in Leipzig, Copenhagen und Braunschweig wurde mit Apparaten, die dem hiesigen ganz gleich sind, beobachtet. Das Resultat der correspondirenden Beobachtungen war dem vom Maitermin angeführten ganz ähnlich. Die zahlreichen in Göttingen beobachteten Schwankungen fanden sich fast alle in den Beobachtungen der andern Plätze wieder, wenn auch in abgeänderten Grössenverhältnissen, doch in unverkennbarer Zusammenstimmung.

Um über dieses merkwürdige Resultat noch ein unabweisbares Zeugniss zu erhalten, wurden bei der damaligen Anwesenheit des Hrn. Prof. \textsc{Weber} in Leipzig einige besondere correspondirende Beobachtungen zwischen Göttingen und Leipzig verabredet, und dazu bestimmte Stunden Vormittags, Mittags und Abends am 1. und 2. October festgesetzt. Diese von vorzüglich eingeübten Beobachtern und mit grösster Sorgfalt ausgeführten Beobachtungen sind in \textsc{Poggendorff}'s Annalen der Physik Bd. 33, S. 426 vollständig abgedruckt, und durch graphische Darstellungen versinnlicht.

Es war hiedurch die Nothwendigkeit, den Erscheinungen in viel engern Zeitintervallen, als Hr. \textsc{von Humboldt} gewählt hatte, zu folgen, auf das klarste vor Augen gelegt. Wir haben eine Zeitlang die Termine in Intervallen von drei zu drei Minuten abgewartet, und dasselbe ist auch von einigen andern Theilnehmern geschehen: da jedoch ein Theil der auswärtigen Theilnehmer sich an die Intervalle von fünf zu fünf Minuten hielt, die auch in den meisten gewöhnlichen Fällen zureichen, so haben wir später der Gleichförmigkeit wegen diese zur allgemeinen Regel angenommen. Da nun aber bei so kleinen Zeitintervallen die
\SourceLeaf{355}Abhaltung der Termine, besonders da, wo nur eine kleine Anzahl von Personen sich in die Arbeit theilen muss, ohne Vergleich mühsamer wird, als beim Aufzeichnen von Stunde zu Stunde, so musste. um das Bestehen des Vereins zu sichern, sowohl die Anzahl als die Dauer der Termine vermindert werden. Die Anzahl ist seit jener Zeit auf sechs im Jahre, die Dauer eines jeden auf 24 Stunden festgesetzt; jedem solchen Haupttermine wurden noch zwei Nebentermine beigefügt. Das Nähere findet man weiter unten.

Nach dieser Einrichtung sind und werden die Beobachtungen ununterbrochen fortgesetzt, in Göttingen und einer fortwährend sich vergrössernden Anzahl anderer Oerter. Apparate, dem Göttingischen gleich oder ähnlich, befinden sich in Altona, Augsburg, Berlin, Bonn, Braunschweig, Breda, Breslau, Cassel, Copenhagen, Dublin, Freiberg, Göttingen, Greenwich, Halle, Kasan, Krakau, Leipzig, Mailand, Marburg, München, Neapel, Petersburg und Upsala. Von acht Oertern aus dieser Anzahl sind bisher noch keine Beobachtungen zu unsrer Kenntniss gekommen, und an einigen der übrigen ist die Theilnahme an den Beobachtungen wegen äusserer Umstände noch keine ununterbrochen regelmässige geworden.

Aus der ersten Zeit dieser Vereinigung sind einige Termine durch \textsc{Schumacher}'s astronomische Nachrichten und \textsc{Poggendorff}'s Annalen der Physik in graphischen Darstellungen veröffentlicht. Seitdem nun aber die Theilnahme sich bereits so sehr vergrössert hat, schien es an der Zeit, auf eine regelmässige Bekanntmachung Bedacht zu nehmen, um die reiche Summe von fruchtbaren Thatsachen zu einem Gemeingut desjenigen Theils des Publicums zu machen, welches sich für die Naturforschung interessirt. Was wir gegenwärtig geben, kann als der erste Jahrgang. seitdem der Verein zu einem gewissen Umfang gekommen ist, betrachtet werden. Vom Jahr 1837 an werden aber die Resultate jedes Termins regelmässig und so bald sie in hinreichender Vollständigkeit beisammen sind, zur Publication gebracht werden.

Die Beobachtungen und ihre graphischen Darstellungen sollen nicht bloss mit denjenigen Erläuterungen und Bemerkungen, welche in einer unmittelbaren Beziehung auf dieselben stehen, begleitet werden, sondern zugleich mit andern Aufsätzen, in welchen Gegenstände aus dem weiten Gebiete des Erdmagnetismus, die darauf bezüglichen Instrumente, ihre Berichtigung und Behandlung, und. die mannigfachen davon zu machenden Anwendungen Platz finden werden.

\SourceLeaf{356}In Beziehung auf den nächsten Gegenstand der Arbeiten unsers Vereins, die Veränderungen in der magnetischen Declination, sei es erlaubt, noch eine Bemerkung hinzuzufügen. Wenn, wie nicht zu bezweifeln ist, die beiden andern Elemente der erdmagnetischen Kraft, die Inclination und die Intensität, ähnlichen Veränderungen unterworfen sind, so kann man fragen, warum vorzugsweise oder für jetzt ausschliesslich, jenem ersten Elemente so sorgfältige Bemühungen gewidmet werden?

Die Kenntniss der Veränderungen und Störungen der magnetischen Declination hat in der That ein sehr grosses praktisches Interesse. Dem Seefahrer, dem Geodäten und dem Markscheider muss ungemein viel daran gelegen sein, zu wissen, wie häufigen und wie grossen Störungen ein Haupthülfsmittel bei seinen Geschäften unvermeidlich unterworfen ist, wäre es auch nur, um das Maass des Vertrauens zu erhalten, welches er demselben schenken darf. Für die beiden letzten Anwendungen der Boussole, in der praktischen Geometrie auf und unter der Erde, kann sogar in Zukunft der Nutzen dieser Untersuchungen noch viel weiter gehen. Wird einmal festgestellt sein, dass die in der Zeit wechselnden unregelmässigen Störungen nie oder nur höchst selten bloss örtlich sind, sondern immer oder fast immer sich in weiten Strecken ganz gleichzeitig und in fast gleicher Grösse offenbaren, so ist das Mittel gegeben, sie fast vollkommen unschädlich zu machen. Der Geodät und der Markscheider braucht nur alle seine Operationen mit der Boussole genau nach der Uhr zu machen und gleichzeitige Beobachtungen an einem andern nicht gar zu entfernten Orte anstellen zu lassen, durch deren Vergleichung jene Störungen sich eben so werden eliminiren lassen, wie reisende Beobachter ihre barometrischen Höhenbestimmungen durch Vergleichung mit Barometerbeobachtungen an einem festen Orte von der unregelmässigen Veränderlichkeit des Barometerstandes unabhängig machen. Dass hier nicht von solchen Störungen die Rede ist, welche die Boussole in eisenhaltigen Gruben erleidet, versteht sich von selbst.

Gleichwohl hat man den Grund des der Declination vor den andern Elementen des Erdmagnetismus gegebenen Vorzuges nicht so wohl in diesen materiellen Rücksichten zu suchen, als vielmehr in dem gegenwärtigen Zustande der Hülfsmittel. Das Aufsuchen der Gesetze in den Naturerscheinungen hat für den Naturforscher seinen Zweck und seinen Werth schon in sich selbst, und ein eigenthümlicher Zauber umgibt das Erkennen von Maass und Harmonie in anschei\SourceLeaf{357}nend ganz Regellosen. Bei der Verfolgung des wunderbaren Spiels in den stets wechselnden Veränderungen der Declination lassen die jetzt angewandten Apparate für Sicherheit, Schärfe und Leichtigkeit der Beobachtungen nichts zu wünschen übrig: allein von den bisherigen Beobachtungsmitteln für die beiden andern Elemente kann man nicht dasselbe sagen. Zur Zeit ist es daher noch zu früh, die letzteren in den Kreis ausgedehnter Untersuchungen aufzunehmen. Sobald aber die Beobachtungsmittel soweit vervollkommnet sein werden, dass wir die Veränderungen, und namentlich die schnell wechselnden Veränderungen, in den andern Elementen des Erdmagnetismus mit Sicherheit erkennen, mit Leichtigkeit verfolgen, und mit Schärfe messen können, werden diese Veränderungen dieselben Ansprüche auf die vereinte Thätigkeit der Naturforscher haben, wie die Veränderungen der Declination. Man darf hoffen, dass dieser Zeitpunkt nicht gar entfernt mehr sein wird.

\clearpage
\fancyhead[L]{\small EIN NEUES HÜLFSMITTEL.}
\SourceLeaf{358}
\pdfbookmark[0]{Göttingische gelehrte Anzeigen. 1837 October 30.}{anzeige-1837-10-30}
\begin{center}
Göttingische gelehrte Anzeigen. 1837 October 30.
\end{center}

Die in der Sitzung der Königl. Societät am 19. September von dem Hofr. \textsc{Gauss} gehaltene Vorlesung hat zum Gegenstande
\begin{center}
\textit{ein neues Hülfsmittel für die magnetischen Beobachtungen.}
\end{center}

Die magnetische Declination, als eines der Elemente der Aeusserungen des Erdmagnetismus, hat nicht allein am frühesten die Beobachter beschäftigt, sondern sie ist, seitdem auch den andern Elementen die Aufmerksamkeit der Naturforscher zugewandt ist, doch in mehreren Beziehungen vor denselben bevorzugt geblieben; von einer der interessantesten Untersuchungen im Gebiete des Erdmagnetismus, über die wunderbaren unregelmässigen, aber über einen ganzen Welttheil gleichzeitig und gleichmässig wirkenden Störungen jener Kraft, wozu in den letzten Jahren ein eigener Verein von Beobachtern zusammen getreten ist, sind die beiden andern Elemente bisher noch ganz ausgeschlossen gewesen. Den Grund dieses der Declination vor den andern Elementen gegebenen Vorzugs hat man weniger in der vielfachen praktischen Wichtigkeit der Kenntniss der Declination für Seefahrer, Geodäten und Markscheider zu suchen, als in dem bisherigen Zustande der Beobachtungsmittel, die, während sie für die Declination Nichts zu wünschen übrig lassen, in Beziehung auf die andern Elemente noch viel weiter zurück sind. Zwar dient das seit einigen Jahren eingeführte Magnetometer, neben seiner Anwendung auf die Bestimmung der Declination, zugleich zur Aus\SourceLeaf{359}mittlung der horizontalen Intensität, und gerade das Problem, diese auf absolutes Maass zurück zu führen, hat zuerst jenes Instrument veranlasst. Allein es löst das Problem noch nicht in \textit{jeder} Beziehung; es kann seiner Natur nach nur einen Mittelwerth der Intensität während eines gewissen Zeitraumes mit Schärfe angeben, und obgleich dieser Zeitraum. gewissermassen von unserer Willkühr abhängt, so darf er doch nicht zu klein genommen werden, weil sonst mit ihm auch die Schärfe und Zuverlässigkeit des Resultats verändert werden würde. Auf die Verfolgung der Veränderungen der Intensität während kurzer Zeitfristen ist daher das Magnetometer gar nicht anzuwenden. Der neue Apparat, welchen die Vorlesung zum Gegenstande hat, ist bestimmt, diese Lücke für die Intensität auszufüllen, und den Beobachtungen dieses Elements dieselbe Leichtigkeit und Sicherheit zu geben, die das Magnetometer für die Declination darbietet.

Der beschränkte Raum dieser Blätter verstattet nicht, hier eine Beschreibung dieses Apparats, den der Hofr. \textsc{Gauss} hat ausführen lassen, zu geben: auch ist dies um so weniger nöthig, da die Vorlesung selbst bald im Druck erscheinen wird. Das Instrument ist seit einigen Monaten in der hiesigen Sternwarte aufgehängt, und bereits in zwei magnetischen Terminen sind die Veränderungen der Intensität an demselben jedesmal 24 Stunden hindurch beobachtet. Verbindet man damit die im magnetischen Observatorium gleichzeitig beobachteten Veränderungen der Declination in Einer Zeichnung, so tritt das wunderbare Spiel der magnetischen Störungen auf eine eigenthümliche neue Art sehr anschaulich hervor, und es lässt sich mit Zuversicht erwarten, dass wenn erst auf ähnliche Weise an mehreren weit von einander entlegenen Orten gleichzeitig beobachtet werden wird, wir über die Sitze der Ursachen dieser räthselhaften Erscheinungen bald umfassendere Aufklärungen gewinnen werden. Während die Intensität sich eben so häufigen und eben so beträchtlichen regellosen Störungen unterworfen zeigt, wie die Declination, tritt doch auch bei jener wie bei dieser das Vorhandensein regelmässig wirkender und mit der Tageszeit zusammen hängender Aenderungen hervor, aber, so viel sich aus täglichen Aufzeichnungen zu bestimmten Stunden während eines Monats erkennen lässt, auf etwas andere Weise. Während nemlich die westliche Declination in unseren Gegenden von Vormittags 7 oder 8 Uhr bis eine oder zwei Stunden nach Mittag zunimmt, und dann wieder zurück geht, ist die Intensität in den ersten Vormittagsstunden abnehmend, erreicht aber ihr Minimum schon eine oder zwei Stunden vor dem Mittage, wo die Declination ge\SourceLeaf{360}rade im raschesten Zunehmen begriffen ist. Es bedarf jedoch kaum der Erinnerung, dass diese Regelmässigkeit an einzelnen Tagen durch die unregelmässigen Störungen oft ganz verdunkelt wird, und genauere Bestimmungen erst die Frucht von lange fortgesetzten Beobachtungen sein können.

Die Einrichtung des Apparats verstattet, denselben ausser seiner Hauptbestimmung noch zu vielen ganz verschiedenen Zwecken anzuwenden. Es ist durch ihn die Auflösung eines Problems gegeben, mit dem man sich früher wiederholt, obwohl ohne Erfolg, beschäftigt hat, nemlich die stündlichen und die unregelmässigen Aenderungen der Declination vergrössert darzustellen. So wie der Apparat gegenwärtig angeordnet ist, beträgt die Vergrösserung das Zehnfache, oder eine Veränderung der Declination, die sich am Magnetometer des magnetischen Observatoriums in 50 Scalentheilen zeigt, erscheint hier mit 500 Scalentheilen. Im letzten magnetischen Termine (30. Sept.) hat man dies durch 8 Stunden gleichzeitiger Beobachtungen an beiden Apparaten auf das befriedigendste bewährt gefunden.

Wie das Magnetometer in Verbindung mit einem Multiplicator bekanntlich ein sehr empfindliches Galvanometer abgibt, eben so der neue Apparat: aber die Empfindlichkeit des letztern in dieser Beziehung übertrifft die des Magnetometers gerade in demselben Verhältniss, wie wir in Beziehung auf Declinationsveränderungen angegeben haben. Der neue Apparat dient also zur scharfen Messung selbst der schwächsten galvanischen Ströme, und es pflegt Bewunderung zu erregen, wie diese den in jenem befindlichen 25 pfündigen Magnetstab in so bedeutende Bewegungen versetzen. In Beziehung auf thermogalvanisch erregte Ströme widerlegt sich dadurch auf das evidenteste die irrige Meinung vieler Physiker, als ob jene eine Kette von bedeutender Länge nicht durchdringen könnten. Durch eine noch so lange Kette werden solche Ströme nicht aufgehoben, sondern nur, und zwar genau, in demselben Verhältnisse geschwächt, wie bei andern Erregungsarten. Unter Anwendung eines thermogalvanischen Apparats von eigenthümlicher Construction bringt die blosse Berührung der Verbindungsstellen mit dem Finger einen galvanischen Strom hervor, der, selbst wenn er eine fast zwei Meilen lange Kette meistens sehr dünnen Drahts zu durchlaufen hat, doch noch in sehr bedeutenden Ablenkungen des Magnetstabes sich zu erkennen gibt.

Die electromagnetischen Wirkungen der gewöhnlichen Reibungselectricität, wenn man sie durch den Multiplicator gehen lässt, gehören zu den schwächsten, schwer zu erkennenden und noch schwerer zu messenden. Bekanntlich ist das
\SourceLeaf{361}Dasein solcher Wirkungen zuerst von \textsc{Colladon} entdeckt und später. von \textsc{Faraday} bestätigt. Anstatt wie diese Physiker gethan haben, eine starke electrische Batterie durch den Leitungsdraht zu entladen, beobachtete der Hofrath \textsc{Gauss} die Wirkung der Reibungselectricität bei fortgesetzter Drehung einer im physicalischen Cabinette aufgestellten Electrisirmaschine, deren Conductor und Reibzeug mit den Enden der grossen, nach der Sternwarte gehenden, Kette verbunden waren. In dieser Kette befand sich der Multiplicator, welcher den Magnetstab des neuen Apparats umgibt, und dieser Stab wurde dadurch in einer Ablenkung von 144 Scalentheilen oder 51 Minuten erhalten, positiver oder negativer, je nach der Richtung, in welcher die Electricität den Multiplicator durchlief. Die Drahtlänge der Kette betrug hierbei etwa 13000 Fuss; aber als besonders merkwürdig muss noch der Umstand hervorgehoben werden, dass eine Verlängerung der Kette bis fast zu einer ganzen Meile, durch Hineinbringen andern Drahts, gar keine Verminderung der electromagnetischen Wirkung hervorbrachte. In dieser Beziehung verhält sich also die strömende Maschinenelectricität anders, als die galvanischen Ströme, die hydrogalvanisch, thermogalvanisch, oder durch Induction erregt werden, und deren durch die magnetische Wirkung gemessene Intensität immer desto schwächer wird, je länger die schliessende Kette ist. Allein weit entfernt, einen wesentlichen inneren Unterschied zwischen jenen und diesen Strömen zu beweisen, dient jene Erscheinung vielmehr zu einer Bestätigung der Gleichheit, und derjenigen Theorie, welcher zufolge ungleiche Intensität zweier Ströme nichts weiter ist, als ungleiche Menge in gegebener Zeit jeden Querschnitt der Leitung durchströmender Electricität. Nur setzen gegebene electromotorische Kräfte der zuletzt genannten Arten desto weniger Electricität in Bewegung, je grösser der Widerstand ist, den die längere Kette entgegensetzt. Aber bei dem oben angeführten Versuche musste \textit{alle} von der Maschine auf den Conductor in Funkenform überspringende Electricität die ganze Kette durchlaufen, um sich mit der entgegengesetzten des Reibzeugs auszugleichen, die Kette mochte kurz oder lang sein (in so fern sie nur hinlänglich isolirt war). Die Menge der in bestimmter Zeit jeden Querschnitt des Leitungsdrahts durchströmenden Electricität hing also gar nicht von der Länge der Kette, sondern nur von dem Spiele der Maschine ab.

Bei den meisten der hier erwähnten Versuche hatte der galvanische oder electrische Strom die grosse zwischen der Sternwarte und dem physicalischen Ca%
